\section*{历史注记}
\phantomsection
\addcontentsline{toc}{section}{历史注记}

（方括号中的数字指本注记末尾的参考文献。）

加法群 \(R^{n}\) 的子群与商群理论的主要结果，自上世纪末以来便已为人所知。算术与分析中的许多问题促使数学家研究由有限多个点生成的 \(R^{n}\) 的子群的结构。例如，拉格朗日在发展“连分数”理论时，顺便证明了：对每个实数 \(\theta\) ，存在不全为零的整数 m、n，使得 \(m - n\theta\) 任意小（{[}1{]}，第~7 卷，第~27 页）。1835 年，雅可比出于他对多复变量周期解析函数的研究，证明了：若 x、y、z 是 \(R^{2}\) 的三个向量，则存在不全为零的整数 m、n、p，使得向量 \(mx + ny + pz\) 任意小（{[}2{]}，第~2 卷，第~25 页）。稍后，狄利克雷在从事代数数理论的工作时，发现了他著名的“抽屉原理”（Schubfachprinzip）（参看 § I，习题 10 与 11），并借助它证明：p 个型 \(\alpha_{i1}m_{1} + \alpha_{i2}m_{2} + \cdots + \alpha_{in}m_{n} - q_{i}\) （ \(1 \leqslant i \leqslant p\) ），其中诸 \(\alpha_{ij}\) 是任意实数，而诸 \(m_{j}\) 与诸 \(q_{i}\) 是（不全为零的）整数，可以同时变得任意小（{[}3{]}，第~I 卷，第~635 页）。埃尔米特于 1850 年用完全不同的方法，对特殊类型的型 \(m\theta_{i} - q_{i}\) （ \(1 \leqslant i \leqslant p\) ）得到了同一结果（{[}4{]}，第~I 卷，第~105 页）。最后，克罗内克于 1884 年证明了 § I 命题 7 所述的一般结果（{[}5{]}，第~31 卷，第~47 页）。

当然，这些结果是独立于局部紧交换群的一般理论而得到的，后者是晚近才出现的（见第三章的历史注记）；但后一理论，特别是对偶理论（*），为这些旧结果带来了新的光明，主要是通过提炼出相伴子群这一基本概念。正文中的论述即基于这些思想（*）。

我们在本章中采取的观点是纯定性的：也就是说，我们确立了 p 个点的带整系数线性组合的存在性，它们可以任意逼近一个给定的点（这个点可能还须满足某些条件）；但我们也可以问，在逼近的精度与给出逼近的线性组合中系数的大小之间是否存在关系；这就是“丢番图逼近”的定量理论的观点，也是第一段中所引的所有作者的观点。在最近一百年里，这些问题成为大量各式各样研究的对象，它们在数论中应用丰富；在此脉络中追溯其发展会使我们远远超出目前的范围，因此，对于希望深入研究这些理论的读者，我们只能向他们指出闵可夫斯基{[}6{]}与 H. 外尔{[}7{]}的基本著作，它们是大量文献的源头（ \(^{**}\) ）。

